
\chapter{机器学习：学习如何学习？}\label{ch:机器学习}

\begin{introduction}
  \item 机器学习的核心概念与数学基础
  \item 从经验中学习的算法原理
  \item 监督学习、无监督学习与强化学习
  \item 优化理论在机器学习中的应用
  \item 泛化能力与过拟合问题的数学分析
\end{introduction}

\begin{introduction}[\faStar\; 你将收获\;\faStar]
  \item 
  \item 
  \item 
  \item 
  \end{introduction}

A computer program is said to {\em learn} from experience~$E$ with respect to some class of tasks~$T$ and performance measure~$P$, if its performance at tasks in~$T$, as measured by~$P$, improves with experience~$E$.  
这句话是机器学习最经典的定义之一。  
它告诉我们：当一个程序通过经验（Experience）不断改进在特定任务（Task）上的表现（Performance），我们就可以说它“学会了”。  
换句话说，学习的本质是——**从经验中提取模式，用以提升未来的表现。**

机器学习，简单来说，就是研究那些能通过经验自动改进的算法。  
对于一个程序，当它面对某个任务~$T$，并用一个标准~$P$来衡量表现时，  
如果随着经验~$E$的积累，它在~$P$下的表现得到提升——我们就说，它从~$E$中“学到了东西”。  
这一定义揭示了机器学习的核心思想：算法不再是死板执行，而是能“越做越好”。

Tom M. Mitchell 在 1997 年的著作中提出了这一定义。  
二十多年过去，它仍被广泛引用，因为它抓住了“学习”的根本——  
即让计算机能够从数据与经验中提取规律，用以改进未来的判断与行为。  
这也正是我们在本章要探讨的核心：  
机器，究竟如何“学习”？又是靠什么实现“改进”的？

在传统编程中，程序员要手动写下每一条规则。  
但在机器学习中，规则不再来自人，而是由算法从数据中自己“悟”出来的。  
这意味着计算机完成了一次质的跨越：  
从“被动执行规则”到“主动生成规则”，  
从遵循人类逻辑到发现数据逻辑。

\begin{figure}[htbp]
  \centering
  
\includegraphics[width=0.9\linewidth]{猫.png}
  \caption{从规则硬编码到数据驱动：图像分类任务中，模型通过数据“学会”识别猫。}
  \label{fig:ml-cat}
\end{figure}

\begin{figure}[htbp]
  \centering
  \includegraphics[width=0.9\textwidth]{手写数字识别.png}
  \caption{手写数字识别示例：机器学习通过样本经验从高维像素中抽取有效特征。}
  \label{fig:ml-mnist}
\end{figure}

这种转变不仅仅是技术层面的突破，更是思维方式的一次革命。  
它迫使我们重新思考：什么才是真正的“学习”？  
这种新观念打破了计算机科学的传统边界，  
将数学、统计学、心理学乃至神经科学都纳入同一个讨论框架。  
%接下来，我们将沿着这一问题展开——机器学习如何定义“学习”，又如何在实践中实现它。
接下来，我们就从整体视角出发，先搭好机器学习的“基本框架”，再逐步展开不同的学习方式与数学工具。
 
\section{从数据到模型：学习的基本框架}

机器学习的核心问题在于：  
我们如何从有限的数据中学习出潜在的规律，  
并利用这些规律去预测未知的情形或未来的变化。  
换句话说，它关注的不是“记住”，而是“概括”——  
如何从已知中推测未知。

在机器学习中，一切都从数据开始。  
数据是学习的原料，也是算法理解世界的窗口。  
一个典型的机器学习系统通常由三部分构成：  
数据集（Dataset）提供经验，模型（Model）表达假设，算法（Algorithm）负责学习与更新。  
三者共同形成了一个“学习闭环”。

数据集为模型提供“经验”；模型决定我们如何表达问题；算法则告诉模型如何从经验中改进。  
它们相互依存，形成了机器学习的基本工作框架。  
理解这一框架，是学习后续各种模型与算法的前提。  
因为几乎所有机器学习方法——无论是监督、无监督还是强化学习——都建立在这三者的互动之上。

从更宏观的角度看，机器学习本质上是一个“函数逼近”问题。  
给定输入 $x$ 与输出 $y$，学习的目标是找到一个函数 $f$，  
使得 $f(x)$ 能够尽可能地接近真实值 $y$。  
换句话说，我们希望 用$f$ 去“拟合”数据背后的规律。  
这就是所谓的学习——找到最能解释数据的函数。

\begin{figure}[htbp]
  \centering
  \includegraphics[width=1\linewidth]{机器学习=一个函数.png}
  \caption{机器学习可被视为寻找从输入到输出的映射函数——以逼近数据背后的规律。}
  \label{fig:ml-as-function}
\end{figure}

\begin{figure}[htbp]
  \centering
  \includegraphics[width=1\textwidth]{LLM函数.png}
  \caption{LLM 亦可视作函数：输入文本、输出概率分布与候选词，体现同一“函数逼近”的思想。}
  \label{fig:llm-function}
\end{figure}

举个例子，在房价预测问题中，$x$代表房屋的特征——  
例如面积、位置、楼层等，而 $y$ 则是对应的房价。  
机器学习要做的，就是通过大量样本 $(x, y)$ 找到一个函数 $f$，  
让它能够根据新的 $x$ 预测出合理的 $y$。  
我们之所以称之为“学习”，  
是因为模型通过经验不断调整自己，使预测越来越接近真实。

在这个学习框架下，算法的任务是调整模型的参数，使预测结果尽量接近真实值。  
为此，我们需要一个衡量“差距”的标准——这就是损失函数（Loss Function）。  
损失函数量化了预测$f(x)$与真实$y$之间的误差，  
而学习的过程，就是不断最小化这种误差的过程。  
可以说，优化思想是机器学习的“心跳”。

从这个角度看，“数据—模型—算法”共同构成了一个学习闭环：  
数据提供经验，模型给出对世界的假设，而算法则负责利用数据来改进这个假设。  
图~\ref{fig:construct-function} 概括了这样一种“构造函数”的过程。  
至于在这个闭环中，如何通过损失函数刻画“好坏”，又如何利用优化算法去改进模型，我们将在后面的“优化理论与损失函数”一节中专门讨论。

损失函数的选择取决于任务类型。  
在回归问题中，我们通常使用均方误差（MSE）来惩罚预测偏差；  
而在分类任务中，交叉熵（Cross-Entropy）则更能反映概率分布之间的差异。  
不同的损失函数意味着不同的学习目标——  
它们引导算法关注不同的方面，从而塑造不同的学习行为。

在优化算法（如梯度下降）的驱动下，  
模型会不断调整自身的参数，让损失函数的值逐步下降，  
直到达到一个相对最优的状态。  
这一过程，就像机器在一次次试探与修正中，逐渐找到通往目标的路径。  
接下来，我们将进一步探讨——这些优化算法究竟是如何实现“学习”的。

\begin{figure}[htbp]
  \centering
  \includegraphics[width=1\textwidth]{机器学习-构造一个函数.png}
  \caption{构造一个函数：数据、模型与优化共同形成学习闭环。}
  \label{fig:construct-function}
\end{figure}

在搭好这张“整体框架图”之后，我们接下来要问的就是：在这套框架之下，根据手头的数据长什么样，具体可以发展出哪些典型的“学习方式”。

\section{监督学习与无监督学习}

从学习所依赖的数据类型来看，  
机器学习最基本的划分方式，是监督学习（Supervised Learning）与无监督学习（Unsupervised Learning）。  
这种区分的核心在于：模型是否依赖带标签的数据进行学习。
换言之，在刚才的“数据—模型—算法”框架之上，我们先根据“数据有没有标签”来区分两种典型情形。

\begin{figure}[htbp]
  \centering
  \includegraphics[width=1\textwidth]{监督与无监督.png}
  \caption{监督 vs 无监督：是否依赖标签是两种学习范式的根本差异。}
  \label{fig:supervised-vs-unsupervised}
\end{figure}

监督学习是最常见、也最直观的一种学习方式。  
它依赖于大量带标签的数据样本，每个样本都包含输入 $x$ 和对应的输出 $y$。  
算法的目标，就是学习从 $x$ 到 $y$ 的映射关系。  
这里的“监督”，指的是标签 $y$ 在训练过程中起到的“指导作用”，  
就像老师在批改作业一样，告诉模型哪里做对了、哪里还需改进。

在监督学习中，最常见的两类任务是分类（Classification）和回归（Regression）。  
分类任务的目标是预测离散的类别标签——  
比如区分“猫”还是“狗”；  
而回归任务则输出连续数值——  
例如预测房价或气温的变化。  
前者回答“是哪一类”，后者回答“有多少”。

举个例子，在垃圾邮件识别任务中，  
我们拥有大量带标签的邮件样本——每封邮件都被标注为“垃圾邮件”或“正常邮件”。  
算法通过学习这些标注样本，建立起从邮件内容到类别的映射关系。  
当新的邮件到来时，它便能自动判断其所属类别。  
这就是监督学习的典型过程：**从已知中学习，以预测未知。**

与监督学习不同，无监督学习并不依赖标签或外部指导。  
它的目标，是在“无人指引”的情况下，从数据中自动发现潜在的结构与规律。  
换句话说，算法在这里不再“模仿答案”，而是自己去探索数据的内在秩序。

无监督学习的典型任务包括聚类（Clustering）与降维（Dimensionality Reduction）。  
聚类通过衡量样本间的相似性，将数据自动分成若干组；  
而降维则试图将高维数据压缩到低维空间，提取其中最具代表性的特征。  
二者虽方法不同，但目标一致——  
都是为了让数据的结构更清晰、更可理解。

以客户分群分析为例，我们往往并没有现成的“标签”。  
算法会根据客户的行为特征和消费记录，自动将他们分为若干群体——  
如“价格敏感型”、“品牌忠诚型”等。  
这类分析帮助企业理解客户结构，从而做出更有针对性的决策。  
这，就是无监督学习在现实世界中的力量。

无监督学习的意义在于，它让机器具备了“自我发现”的能力。  
它揭示了数据中人类未曾标注、却真实存在的结构与模式，  
为后续的预测、推荐与决策提供了深层的知识基础。  
从这个角度看，无监督学习不仅是一种技术，  
更是一种“让数据自己说话”的方法论。

监督学习与无监督学习，主要处理的是“静态的数据集”：数据先收集好，再交给算法进行学习。  
在很多场景中，这已经非常有用。  
但在另一些情形下，智能体需要一边行动、一边获得反馈，在与环境不断互动中学习决策策略——这就把我们自然引向了下一种学习方式：强化学习。

\section{强化学习：从试错中学习}

强化学习（Reinforcement Learning, RL）是一种以“试错”为核心机制的学习方法。  
它让智能体在与环境的不断交互中，逐步学会在不同情境下采取能带来最大奖励的行动。  
之所以称为“强化”，是因为有益的行为会被奖励并得到强化，而无效的行为则会被自然淘汰。

在强化学习框架中，学习的主体称为“智能体”（Agent），  
它所面对的外部世界称为“环境”（Environment）。  
智能体在环境中感知状态（State），选择行动（Action），并根据反馈信号（Reward）进行调整。  
这种循环过程，正如人类在不断尝试与反馈中逐渐学会一项技能——  
它的本质，就是从经验中改进行为。

强化学习的核心目标，是学习到一个最优的策略（Policy）。  
所谓“策略”，就是智能体在每种状态下应采取什么行动的规则。  
学习的过程，实际上就是寻找一个能在长期交互中获得最高累计奖励的决策方式。

形式上，强化学习通常建模为一个马尔可夫决策过程（Markov Decision Process, MDP）。  
它由四个要素构成：状态空间、动作空间、状态转移概率以及奖励函数。  
采用 MDP 的原因在于，它能够刻画“当前决策只依赖当前状态”的情境，  
这为算法提供了清晰的数学框架。

在每个时间步 $t$，智能体观察到状态 $s_t$，执行动作 $a_t$。  
环境则根据转移概率 $P(s_{t1}|s_t, a_t)$ 变化到新状态 $s_{t1}$，并返回奖励 $r_t$。  
智能体的目标，是找到一种策略 $\pi$，使得未来所有奖励的折扣和最大化：
$$
\max_\pi \mathbb{E}\left[\sum_{t=0}^\infty \gamma^t r_t\right]
$$  
其中，$\gamma$ 是折扣因子（$0 < \gamma < 1$），  
用于平衡“眼前的收益”和“长远的回报”。

强化学习的关键挑战在于平衡“探索”（Exploration）与“利用”（Exploitation）。  
智能体必须敢于尝试新行动，以发现潜在的更优策略；  
同时，又要善于利用已有经验，以获得稳定的回报。  
这就像人类学习：一味重复旧路无法进步，而盲目尝试又可能走弯路。

强化学习的研究发展出了多种算法框架。  
其中，值函数法（Value-based methods）侧重于评估每个状态或动作的价值；  
策略梯度法（Policy Gradient）则直接优化策略本身；  
而演员-评论家法（Actor-Critic）融合了二者的思想——  
由“演员”负责行动，“评论家”负责评估，形成一个相互促进的学习体系。

强化学习已在多个领域展现出惊人的成果。  
从击败人类顶级棋手的 AlphaGo，到实现动态决策的自动驾驶系统，再到精确操控的智能机器人——  
这些都是强化学习的成功应用。  
它们的共同点在于：环境复杂、反馈往往延迟、目标依赖长期规划。  
而强化学习，正擅长在这种“不确定而连续”的情境中成长。

强化学习的重要性，不止于解决复杂任务。  
更深层的意义在于——它揭示了智能的形成过程：  
学习不是静态的积累，而是与环境的动态互动。  
在不断试探、反馈与调整中，智能体学会了决策、形成了经验，  
这也为我们理解“智能”本身，提供了一个可计算的模型。

无论是监督学习、无监督学习，还是强化学习，  
它们背后都有一个共同的问题：  
如何用一个合适的“目标函数”来衡量好坏，并通过优化算法不断改进模型。  
接下来，我们就从“优化”的角度，再看一看机器学习中的学习过程。

\section{优化理论与损失函数}

优化，是机器学习的核心思想之一。  
几乎所有学习算法，都可以理解为——在某个目标函数之下，不断寻找最优解的过程。  
换句话说，机器学习的本质，就是通过数学方法去“逼近最好”。

这个目标函数，通常被称为损失函数（Loss Function）或代价函数（Cost Function）。  
它衡量的是模型预测值与真实结果之间的差距。  
损失越小，说明模型的预测越接近现实，也就意味着学习越成功。

优化的目标，就是找到一组参数 $\theta$，  
使损失函数 $L(\theta)$ 的值最小化：
$$
\min_{\theta} L(\theta)
$$  
这里的 $\theta$ 代表模型的可调部分——  
比如神经网络的权重、线性回归的系数等。  
优化，就是调整这些参数，让模型在数据上表现得更好。

在众多优化方法中，梯度下降法（Gradient Descent）是最基础也最常用的。  
它通过计算损失函数关于参数的梯度，  
沿着梯度的反方向逐步更新参数。  
这就像一个登山者在迷雾中摸索下山——  
每一步都沿着“坡度”下降，直到抵达谷底。

梯度下降的更新规则如下：
$$
\theta_{t1} = \theta_t - \eta \nabla_\theta L(\theta_t)
$$  
其中 $\eta$ 称为学习率（Learning Rate），它决定每次参数更新的步长。  
学习率太大，模型可能“跳过”最优点；  
太小，又会导致学习速度过慢。  
因此，选择合适的学习率，是优化的艺术之一。

当损失函数的形状复杂时，  
梯度下降法有时会陷入局部最小值，或在狭长的谷地中反复震荡。  
为此，研究者们设计了许多改进算法，  
如动量法（Momentum）、Adam、RMSprop 等。  
它们在不同程度上解决了“速度”与“稳定性”的平衡问题，  
让优化过程更加高效、平滑。

动量法（Momentum）在更新过程中引入了“惯性”概念。  
它不仅考虑当前梯度，还结合了前一次更新的方向：  
$$
v_{t1} = \beta v_t  (1 - \beta) \nabla_\theta L(\theta_t)
$$  
$$
\theta_{t1} = \theta_t - \eta v_{t1}
$$  
这种机制让参数像带有“动量”的小球，在下降过程中能跨越小陷阱，  
减少震荡，更快接近全局最优。

Adam 算法可以看作是动量法与自适应学习率的结合体。  
它会根据每个参数的梯度历史，自主调整学习率的大小，  
从而在不同任务和数据分布下，都能保持较高的稳定性与收敛速度。  
正因如此，Adam 几乎成为现代深度学习的“默认选择”。

除了优化算法本身，损失函数的设计也至关重要。  
它决定了“什么才算学得好”。  
不同的任务，需要不同形式的损失函数，  
而如何设计它，往往体现着我们对问题本质的理解。

例如，在回归问题中，我们常使用均方误差（MSE）：
$$
L = \frac{1}{n}\sum_{i=1}^{n}(y_i - \hat{y}_i)^2，
$$  
它强调预测值与真实值之间的距离。  
而在分类问题中，交叉熵（Cross-Entropy）更常见：
$$
L = -\sum_{i} y_i \log(\hat{y}_i)
$$  
它从信息论的角度，衡量两个概率分布之间的差异。  
两者的选择，体现了任务类型与目标本质的不同。

损失函数与优化算法，共同构成了机器学习的“学习引擎”。  
前者决定方向，后者决定速度。  
一个恰当的损失函数能让模型朝着正确目标前进，  
而一个高效的优化方法，则让它更快、更稳地抵达终点。  
%两者协同，才使学习真正发生。
当然，把训练集上的损失降到很低，只是“学会了教材上的例题”；  
要判断模型是否真正学会了任务本身，还需要考察它在“新题”上的表现——这就引出了泛化与过拟合的问题。

\section{泛化能力与过拟合}

机器学习的真正目标，并非让模型在训练数据上“背熟题”，  
而是希望它能在未见过的新样本上，也保持良好的表现。  
这种能力被称为“泛化能力”（Generalization）。  
泛化，是从经验中“触类旁通”的能力，  
也是判断一个模型是否真正学会的关键标准。

如果一个模型在训练集上表现极好，  
却在新数据上显著失效，我们称这种现象为“过拟合”（Overfitting）。  
这就像学生死记硬背了题目答案，考试时题一变就不会了——  
模型“记住”了数据，却没有真正“理解”规律。

与之相对的，是“欠拟合”（Underfitting）。  
当模型在训练集和测试集上都表现不佳时，  
往往意味着模型太简单，无法捕捉数据背后的模式。  
换句话说，它还没有真正“学会”任务本身。

过拟合与欠拟合，揭示了模型复杂度与泛化能力之间的平衡关系。  
模型太复杂，容易被噪声牵着走；太简单，又难以理解数据结构。  
理想的模型，应该既能学习规律，又能忽略偶然。  
这正是模型设计与调优的艺术所在。

\begin{figure}[htbp]
  \centering
  \includegraphics[width=0.9\textwidth]{过(欠)拟合.png}
  \caption{欠拟合与过拟合：模型复杂度与泛化能力的平衡。}
  \label{fig:under-over-fitting}
\end{figure}

为了提升模型的泛化能力，我们常引入正则化（Regularization）技术。  
它通过在损失函数中加入对模型复杂度的惩罚项，  
抑制参数过度拟合数据中的噪声：
$$
L'(\theta) = L(\theta)  \lambda \|\theta\|^2
$$  
这种做法看似“让模型退一步”，  
却往往能让它在新数据上“进两步”。

式中的 $\lambda$ 称为正则化系数，用来控制惩罚项的强度。  
$\lambda$ 过大，模型被“约束”得太紧，容易欠拟合；  
$\lambda$ 过小，则放任参数自由，容易过拟合。  
因此，找到合适的 $\lambda$，  
就是在“自由”与“约束”之间取得平衡的过程。

正则化有多种形式，其中最常见的是~$L_1$ 正则化（Lasso）与~$L_2$ 正则化（Ridge）。  
~$L_1$正则化会让部分参数趋近于零，从而实现“自动特征选择”；  
而~$L_2$ 正则化则让参数整体更平滑，避免出现过大的权重。  
从直觉上看，前者像“取舍”，后者像“均衡”。

除了正则化，  
我们还可以通过交叉验证（Cross Validation）评估模型稳健性，  
通过早停（Early Stopping）防止过度训练，  
或通过数据增强（Data Augmentation）扩展样本多样性。  
这些方法的共同目标，都是让模型学得“恰到好处”——  
既不过度依赖训练数据，又能适应现实变化。

泛化能力，是机器学习的灵魂。  
一个真正智能的系统，不应只会记忆，更要学会举一反三。  
只有具备良好泛化能力的模型，  
才能称得上“从经验中学习”，而不仅仅是“重演经验”。  
这，正是机器学习区别于简单数据拟合的关键所在。

\section{本章小结：学习的数学本质}

从监督学习到强化学习，  
我们一路看到：无论方法多么不同，本质上都是在追求“最优”的数学过程。  
换句话说，学习的多样性背后，隐藏着一个统一逻辑——  
都是在不同条件下寻找最优解。

监督学习通过最小化误差，去逼近一个理想的函数；  
无监督学习通过挖掘数据结构，实现自发的特征抽取；  
强化学习则在不断交互中，优化自己的行为策略。  
三者看似路径不同，实则都在更新模型对世界的理解——  
学习，就是在“误差—调整—再试”的循环中前行。

这也意味着，学习远非简单的模式记忆。  
它是一种在高维空间中，通过不断优化，寻找最优点的过程。  
每一次迭代，都是对假设的修正，对模式的再理解。  
这正体现了数学思维的魅力：  
用抽象的逻辑去揭示复杂系统中的秩序。

学习的核心是“优化”，而优化的核心是“反馈”。  
机器通过误差信号更新参数，人类则通过经验修正认知。  
两者的共同点在于——都依赖对结果的反思。  
无论是算法还是人脑，真正的学习，  
都发生在反馈的瞬间。

因此，从数学的视角来看，  
学习可以被理解为一个在不确定性中持续优化决策的过程。  
它让数据转化为知识，让模式转化为预测，  
也让智能系统具备了“自我改进”的能力。  
理解这一点，我们便不再只是“使用”机器学习，  
而是开始真正理解——何为“学习”。

机器学习，不仅是一套算法体系，  
更是一种思维方式。  
它让我们重新审视“学习”这一人类最古老的能力，  
也让数学——这门最理性的语言——在智能时代焕发出新的生命力。  
或许，理解机器如何学习，  
也是理解我们自身如何思考的一面镜子。

\nocite{*}
%\printbibliography[heading=bibintoc, title=\ebibname]
\newpage
